\section{Conclusion}

In this study, we established a comprehensive multi-omics systems biology framework that integrates the NASA OSD-615 (APEX-03-1) high-throughput cell wall glycome dataset with companion spaceflight transcriptomics, Salk Institute single-cell spatial transcriptomics, protein interactome modeling, dynamic vesicle transport simulations, and TabPFN tabular foundation model cross-mission meta-analysis.

Our findings demonstrate that:
\begin{enumerate}
    \item Microgravity profoundly alters plant cell wall polysaccharide deposition, specifically accelerating secondary wall xylan and cellulose deposition while reorganizing primary wall xyloglucan and arabinogalactan epitope extractability.
    \item Single-cell spatial projection onto the Salk Institute root atlas and \textit{ggPlantmap} vector anatomy pinpoints differentiating metaxylem vessels as the primary anatomical epicentre of secondary wall xylan hyper-accumulation, directly confirmed by in situ confocal immunohistochemistry ($+192\%$ xylem xylan increase; $r = +0.87$ with ELISA).
    \item These wall matrix modifications are directly correlated with transcriptional activation of class XI myosins (\textit{MYA1}, \textit{MYA2}), secondary wall synthases (\textit{CESA4}, \textit{IRX9}), and nucleocytoplasmic glycosyltransferases (\textit{SEC}, \textit{SPY}).
    \item Dynamic Monte Carlo simulations show that microgravity-induced cortical microtubule disorientation decreases secretory vesicle delivery efficiency by 28\% and extends transit times by 71\%, establishing a quantitative biophysical link between cytoskeletal track integrity and cell wall remodeling.
    \item TabPFN foundation model cross-study meta-analysis with independent spaceflight study OSD-121 (STS-131 BRIC-16) achieved perfect zero-shot transfer ($\text{ROC-AUC} = 1.000$) and identified a sharp non-linear planetary gravity threshold, proving that Lunar gravity ($0.16g$) triggers near-complete microgravity adaptation, whereas Martian gravity ($0.38g$) preserves normal terrestrial root developmental programs.
    \item We provide an exhaustive analytical roadmap of advanced mass spectrometry workflows (EThcD, lectin chromatography, click-chemistry) and AlphaFold 3D structural insights to facilitate future orbital glycoproteomics studies.
\end{enumerate}

To ensure complete adherence to FAIR (Findable, Accessible, Interoperable, Reusable) open science principles, all curated datasets, scripts, simulation engines, automated test suites, and interactive web visualization tools are publicly deposited on GitHub and Zenodo. This open-science platform will empower the space biology community to design resilient crop ideotypes for bioregenerative life support systems (BLSS) on upcoming Artemis lunar outposts and long-duration Mars transit missions.

\section*{Data Availability}
All raw and processed datasets analyzed in this study are publicly available in the NASA Open Science Data Repository (OSDR) under accessions OSD-615 (\url{https://osdr.nasa.gov/bio/repo/data/studies/OSD-615}), OSD-218 (\url{https://osdr.nasa.gov/bio/repo/data/studies/OSD-218}), OSD-217 (\url{https://osdr.nasa.gov/bio/repo/data/studies/OSD-217}), and OSD-121 (\url{https://osdr.nasa.gov/bio/repo/data/studies/OSD-121}). The Salk Institute single-cell root transcriptomic reference atlas is accessible via the Gene Expression Omnibus (GEO accession GSE243171) and \textit{Nature Plants} (\url{https://doi.org/10.1038/s41477-025-02072-z}). Curated matrices, RO-Crate 1.1 JSON-LD metadata, and pre-computed models are deposited in Zenodo and included in the project repository.

\section*{Code Availability}
All Python analysis scripts, stochastic simulation models, automated pytest verification suites, and web dashboard source code are open-source under the MIT License at \url{https://github.com/dr-richard-barker/OSD615-glycome-cytoskeleton-systems-biology}. The interactive systems biology dashboard is deployed live at \url{https://dr-richard-barker.github.io/OSD615-glycome-cytoskeleton-systems-biology/}.
